\documentclass[10pt,letterpaper]{article}

\usepackage[utf8]{inputenc}
\usepackage[T1]{fontenc}
\usepackage[margin=1in]{geometry}
\usepackage{amsmath,amssymb,amsfonts}
\usepackage{booktabs}
\usepackage{multirow}
\usepackage{graphicx}
\usepackage{xcolor}
\usepackage{hyperref}
\usepackage{algorithm}
\usepackage{algorithmic}
\usepackage{natbib}
\usepackage{times}

\newcommand{\ours}{ACR-VLA}
\newcommand{\crm}{CRM}
\newcommand{\cgs}{CGS}
\newcommand{\aci}{ACI}
\newcommand{\coc}{COC}
\newcommand{\ie}{\textit{i.e.}}
\newcommand{\eg}{\textit{e.g.}}
\newcommand{\etal}{\textit{et al.}}

\title{ACR-VLA: Dimension-Selective Action Correction via\\Per-DOF Conformal Features for Robust Vision-Language-Action Models}

\author{
  Anonymous Authors
}

\begin{document}

\maketitle

% ============================================================
\begin{abstract}
Flow-matching Vision-Language-Action (VLA) models achieve near-perfect performance on standard manipulation benchmarks but suffer catastrophic degradation under distribution shift.
Existing robustness approaches either operate at training time (limiting inference-time adaptability) or reduce rich uncertainty information to scalar summaries (discarding cross-dimensional structure).
We present \ours{}, an inference-time action correction framework that leverages \emph{per-DOF conformal features} extracted from multivariate conformal prediction on flow-matching samples.
A Conformal Geometry Summarizer (\cgs{}) computes 45-dimensional structured uncertainty features---per-DOF standard deviations, conformal radii, and ellipsoid volume---from $K{=}10$ flow-matching samples using Ledoit-Wolf shrinkage estimation.
A lightweight Conformal Refinement Module (\crm{}, 131K parameters) conditions on these features to produce dimension-selective action corrections, aggressively adjusting high-uncertainty DOFs while preserving confident predictions.
Adaptive Conformal Inference (\aci{}) with Chunk Overlap Consistency (\coc{}) maintains calibrated uncertainty online without external supervision.
On LIBERO-Plus, \ours{} improves base VLA success rate by 11.8 percentage points (pp) while degrading nominal performance by only 0.4pp.
When stacked on SRPO---the current training-time state of the art at 82.1\%---\ours{} achieves 86.4\%, demonstrating that inference-time conformal correction is orthogonal to training-time optimization.
Ablations confirm a clear hierarchy: per-DOF features (+7.3pp over no-conformal) $>$ scalar uncertainty (+4.4pp over scalar) $>$ no correction, providing the first evidence that structured multivariate conformal statistics serve as effective learned correction features for robotic action generation.
\end{abstract}

% ============================================================
\section{Introduction}
\label{sec:intro}

Flow-matching Vision-Language-Action (VLA) models have achieved remarkable success on standard robotic manipulation benchmarks, with recent methods reaching 97--99\% success rates on LIBERO~\cite{liu2024libero}.
However, the LIBERO-Plus benchmark~\cite{liangplus2025} exposes a critical robustness gap: even state-of-the-art VLAs degrade to 50--70\% success under realistic distribution shifts including camera perturbations, observation noise, and robot-state offsets.

Existing approaches to VLA robustness fall into two categories, each with fundamental limitations.
\textbf{Training-time methods}---such as SRPO~\cite{SRPO}, RobustVLA~\cite{RobustVLA}, and DiG-Flow~\cite{digflow}---improve robustness through augmented training or adversarial optimization.
While effective, they require retraining and cannot adapt to novel perturbations at deployment.
\textbf{Inference-time methods}---such as ReconVLA~\cite{reconvla} and SAFE~\cite{SAFE}---detect or mitigate distribution shift during execution.
However, ReconVLA reduces uncertainty to a scalar Euclidean distance for sample selection, discarding the rich per-dimension structure of the flow-matching posterior.
SAFE detects anomalies but does not correct actions.

We identify a key gap: no existing method leverages the \emph{structured, per-DOF uncertainty} inherent in flow-matching VLA samples to perform \emph{learned, dimension-selective action correction} at inference time.
Multivariate conformal prediction on flow-matching samples naturally yields per-DOF conformal radii and cross-DOF covariance structure that encode \emph{which action dimensions are uncertain and how they co-vary}---information fundamentally richer than any scalar summary.

We propose \ours{} (\textbf{A}daptive \textbf{C}onformal \textbf{R}efinement for \textbf{VLA}s), the first method to use structured multivariate conformal prediction statistics as input features for a learned action correction module.
Our approach consists of three components:
\begin{enumerate}
    \item \textbf{Conformal Geometry Summarizer (\cgs{})}: Computes 45-dimensional per-DOF uncertainty features from $K{=}10$ flow-matching samples using Ledoit-Wolf shrinkage covariance estimation and Mahalanobis-based conformal prediction.
    \item \textbf{Conformal Refinement Module (\crm{})}: A lightweight 131K-parameter network that conditions on \cgs{} features to produce dimension-selective action residuals.
    \item \textbf{Adaptive Conformal Inference (\aci{}) with Chunk Overlap Consistency (\coc{})}: An online calibration mechanism that maintains valid conformal coverage using self-referential consistency signals, requiring no external supervision.
\end{enumerate}

Our contributions are:
\begin{itemize}
    \item We demonstrate that per-DOF conformal features enable dimension-selective correction, with ablations showing a clear hierarchy: per-DOF (74.1\%) $>$ scalar (69.7\%) $>$ no-conformal (66.8\%) $>$ base (62.3\%) on LIBERO-Plus.
    \item We show that \aci{}+\coc{} maintains calibrated uncertainty online (Spearman $\rho{=}0.62$ between \coc{} and task success; 31\% ellipsoid volume narrowing over static calibration) without external supervision.
    \item We establish that \ours{} is orthogonal to training-time methods, pushing SRPO from 82.1\% to 86.4\% on LIBERO-Plus---a new state of the art.
    \item We provide mechanistic evidence via gradient-based feature importance analysis showing that \crm{} attends to perturbation-specific DOF subsets, confirming emergent dimension-selective behavior.
\end{itemize}

% ============================================================
\section{Related Work}
\label{sec:related}

\paragraph{Robust VLA Training.}
Training-time approaches improve VLA robustness through data augmentation and optimization.
SRPO~\cite{SRPO} uses self-play reinforcement to optimize for perturbation-invariant policies, achieving 82.1\% on LIBERO-Plus.
RobustVLA~\cite{RobustVLA} applies adversarial training with domain-randomized observations.
DiG-Flow~\cite{digflow} integrates gradient-based robustness objectives into flow-matching training.
These methods are complementary to \ours{}; our stacking experiments (Section~\ref{sec:stacking}) confirm orthogonality.

\paragraph{Inference-Time Uncertainty for Robotics.}
ReconVLA~\cite{reconvla} generates multiple flow-matching samples and selects the one with minimum Euclidean distance to the sample mean---a scalar uncertainty measure that discards per-DOF structure.
SAFE~\cite{SAFE} uses reconstruction error for anomaly detection but does not correct actions.
UPS~\cite{UPS} applies conformal prediction for policy gating but does not condition a learned correction on conformal features.
A2C2~\cite{A2C2} and PLD~\cite{PLD} learn residual corrections but without calibrated uncertainty conditioning.
CP-in-the-Loop~\cite{cploop} uses conformal sets for planning feedback but is not a learned parametric correction.

\paragraph{Conformal Prediction in Robotics.}
Conformal prediction provides distribution-free coverage guarantees and has been applied to robot planning~\cite{cploop}, safety monitoring~\cite{conformalsafety}, and trajectory prediction~\cite{conformaltrajectory}.
Most applications use univariate or Bonferroni-corrected marginal CP.
Our work is distinguished by (a) using \emph{multivariate} conformal prediction with joint Mahalanobis statistics, and (b) treating conformal statistics as \emph{input features} for a learned module rather than decision boundaries.

\paragraph{Disentangling Uncertainty Types.}
DiU~\cite{DiU} estimates aleatoric uncertainty from single-pass inference via distributional heads.
Our ablation (Section~\ref{sec:epistemic_aleatoric}) demonstrates that epistemic uncertainty from flow-matching sample diversity is the appropriate signal for distribution-shift correction ($+6.7$pp over aleatoric features), consistent with the theoretical distinction between model uncertainty (reducible with more data) and inherent stochasticity.

% ============================================================
\section{Method}
\label{sec:method}

\subsection{Preliminaries: Flow-Matching VLAs}
\label{sec:prelim}

Flow-matching VLA policies $\pi_\theta$ learn a velocity field $v_\theta(a_t, t \mid o)$ that transports samples from a prior $p_0 = \mathcal{N}(0, I)$ to the action distribution conditioned on observation $o$.
Actions are generated by integrating the ODE $da/dt = v_\theta(a, t \mid o)$ from $t{=}0$ to $t{=}1$.
We generate $K{=}10$ samples $\{a^{(k)}\}_{k=1}^K$ from independent initial noise draws, producing a set of 7-dimensional action vectors (3 position + 3 rotation + 1 gripper).

\subsection{Conformal Geometry Summarizer (CGS)}
\label{sec:cgs}

\cgs{} extracts structured uncertainty features from the $K$ flow-matching samples in three steps.

\paragraph{Robust Covariance Estimation.}
With $K{=}10$ samples in $p{=}7$ dimensions ($p/n{=}0.7$), the sample covariance is poorly conditioned.
We apply Ledoit-Wolf shrinkage~\cite{ledoitwolf}:
\begin{equation}
    \hat{\Sigma}_{\text{LW}} = (1-\delta)\hat{\Sigma} + \delta \cdot \text{diag}(\hat{\Sigma}), \quad \delta = \arg\min_{\delta'} \mathbb{E}\|\hat{\Sigma}_{\delta'} - \Sigma\|_F^2
\end{equation}
where $\hat{\Sigma}$ is the sample covariance and $\delta \in [0,1]$ is the data-driven shrinkage intensity.
Empirically, $\delta \approx 0.38$ on average, indicating moderate regularization that preserves dominant cross-DOF structure (median condition number 14.7).

\paragraph{Multivariate Conformal Prediction.}
Given a calibration set of base-policy rollouts $\{(o_i, a_i^*)\}_{i=1}^n$, we compute nonconformity scores via the Mahalanobis distance:
\begin{equation}
    s_i = (a_i^* - \bar{a}_i)^\top \hat{\Sigma}_{\text{LW},i}^{-1} (a_i^* - \bar{a}_i)
\end{equation}
where $\bar{a}_i$ is the sample mean of $K$ flow-matching samples at calibration point $i$.
The conformal threshold $\hat{q}_\alpha$ is the $\lceil(1-\alpha)(n+1)\rceil/n$ quantile of $\{s_i\}$.

\paragraph{45-Dimensional Feature Extraction.}
From the per-step covariance $\hat{\Sigma}_{\text{LW}}$ and conformal threshold $\hat{q}_\alpha$, we extract:
\begin{itemize}
    \item Per-DOF sample standard deviation: $\sigma_j = \sqrt{[\hat{\Sigma}_{\text{LW}}]_{jj}}$ for $j=1,\ldots,7$ \hfill (7D)
    \item Per-DOF conformal radius: $r_j = \hat{q}_\alpha \cdot \sigma_j \cdot \sqrt{[\hat{\Sigma}_{\text{LW}}^{-1}]_{jj}^{-1}}$ \hfill (7D)
    \item Per-DOF-pair correlation: $\rho_{jk} = [\hat{\Sigma}_{\text{LW}}]_{jk}/(\sigma_j\sigma_k)$ for $j<k$ \hfill (21D)
    \item Eigenvalue spectrum of $\hat{\Sigma}_{\text{LW}}$ \hfill (7D)
    \item Ellipsoid volume: $\hat{q}_\alpha^{p/2} \cdot |\hat{\Sigma}_{\text{LW}}|^{1/2} \cdot V_p$ \hfill (1D)
    \item Log-determinant: $\log|\hat{\Sigma}_{\text{LW}}|$ \hfill (1D)
    \item Condition number: $\lambda_{\max}/\lambda_{\min}$ \hfill (1D)
\end{itemize}
These 45 features encode per-DOF uncertainty magnitudes, cross-DOF covariance structure, and global uncertainty geometry---the information needed for dimension-selective correction.

\subsection{Conformal Refinement Module (CRM)}
\label{sec:crm}

\crm{} is a lightweight feedforward network (131K parameters) that maps the 45D conformal features and base action $\bar{a}$ to an action residual $\Delta a$:
\begin{equation}
    a_{\text{corrected}} = \bar{a} + \crm{}([\bar{a}; f_{\cgs{}}]; \phi)
\end{equation}
where $f_{\cgs{}} \in \mathbb{R}^{45}$ is the \cgs{} feature vector and $\phi$ are learned parameters.
The architecture consists of two hidden layers (256 units each, ReLU activation) with a linear output head producing a 7D residual.

\paragraph{Training.}
\crm{} is trained on LIBERO demonstrations augmented with 5 synthetic perturbation types (camera jitter, Gaussian noise, robot-state offset, lighting shift, background swap).
For each training sample, we (a) apply a random perturbation, (b) generate $K{=}10$ flow-matching samples from the perturbed observation, (c) extract \cgs{} features, and (d) supervise the corrected action against the ground-truth demonstration action:
\begin{equation}
    \mathcal{L}(\phi) = \mathbb{E}_{(o,a^*) \sim \mathcal{D}, \xi \sim \text{Aug}} \left\| a^* - \left(\bar{a}_\xi + \crm{}([\bar{a}_\xi; f_{\cgs{},\xi}]; \phi)\right) \right\|_2^2
\end{equation}

The intentionally small capacity (131K parameters) provides an inductive bottleneck that prevents overfitting and ensures the module learns to rely on conformal features rather than memorizing corrections.

\subsection{Adaptive Conformal Inference with Chunk Overlap Consistency}
\label{sec:aci}

Static conformal calibration becomes miscalibrated under distribution shift.
We apply Adaptive Conformal Inference (\aci{})~\cite{gibbs2021aci} to update the conformal threshold online:
\begin{equation}
    \hat{q}_{\alpha,t+1} = \hat{q}_{\alpha,t} + \gamma(\alpha - \text{err}_t)
\end{equation}
where $\text{err}_t$ is a coverage error signal and $\gamma$ is the learning rate.

\paragraph{Chunk Overlap Consistency (COC).}
Flow-matching VLAs generate action \emph{chunks}---sequences of $H$ future actions.
At overlapping timesteps, consecutive chunks predict the same future actions.
We define the COC coverage signal: a timestep is ``covered'' if the action predicted by the previous chunk falls within the current chunk's conformal set at the overlapping position.
This provides a dense, self-referential coverage signal (${\sim}200$ updates per episode) without requiring ground-truth actions or task success labels.

Empirically, COC coverage correlates with task success (Spearman $\rho{=}0.62$, $p{<}10^{-4}$), validating its use as a proxy for prediction quality.
\aci{} with COC converges in ${\sim}45$ episodes versus ${\sim}120$ with episode-level binary feedback, and achieves 31\% tighter conformal ellipsoids at matched 90\% marginal coverage.

\subsection{Inference Pipeline}
\label{sec:inference}

At each control step:
(1) Generate $K{=}10$ flow-matching samples from the current observation.
(2) Compute Ledoit-Wolf covariance and extract 45D \cgs{} features using the current \aci{}-adapted threshold.
(3) Pass base action (sample mean) and \cgs{} features through \crm{} to produce the corrected action.
(4) Execute the corrected action; update \aci{} threshold via COC signal.

Total inference overhead of \crm{} is 3.2ms (on top of the flow-matching sampling); the full pipeline runs at 385ms per step with 100-step ODE integration.
With SnapFlow 1-step distillation, total inference reduces to 45ms ($8.6\times$ speedup, $-1.6$pp).

% ============================================================
\section{Experiments}
\label{sec:experiments}

\subsection{Setup}
\label{sec:setup}

\paragraph{Benchmarks.}
We evaluate on LIBERO~\cite{liu2024libero} (10 tasks, nominal conditions) and LIBERO-Plus~\cite{liangplus2025} (7 perturbation types at severity levels L3--L5, 5 tasks $\times$ 20 rollouts per condition).
LIBERO-Plus perturbations span camera viewpoint, observation noise, robot-state offset, lighting, background, object layout, and language paraphrase---covering both sensor-level and semantic distribution shifts.

\paragraph{Base Model.}
We use SmolVLA~\cite{smolvla} as the base flow-matching VLA, achieving 97.2\% on nominal LIBERO.
For stacking experiments, we also apply \ours{} to the SRPO-finetuned~\cite{SRPO} SmolVLA checkpoint (98.1\% nominal, 82.1\% LIBERO-Plus).

\paragraph{Baselines.}
We compare against: (a)~SmolVLA base (no correction), (b)~SmolVLA + ReconVLA~\cite{reconvla} (scalar-uncertainty sample selection), and (c)~SRPO~\cite{SRPO} (training-time SOTA on LIBERO-Plus at 82.1\%).

\subsection{Main Results}
\label{sec:main_results}

\begin{table}[t]
\centering
\caption{Main results on LIBERO (nominal) and LIBERO-Plus (distribution shift). SR denotes success rate (\%). \ours{} improves LIBERO-Plus performance by 11.8pp over the base VLA with minimal nominal degradation.}
\label{tab:main}
\begin{tabular}{lccccc}
\toprule
\textbf{Method} & \textbf{LIBERO} & \textbf{L-Plus L3} & \textbf{L-Plus L4} & \textbf{L-Plus L5} & \textbf{L3--L5 Avg} \\
\midrule
SmolVLA (base) & 97.2 & 71.8 & 60.4 & 54.7 & 62.3 \\
SmolVLA + ReconVLA & 96.5 & 75.2 & 64.8 & 59.1 & 66.4 \\
\textbf{SmolVLA + \ours{} (ours)} & \textbf{96.8} & \textbf{82.6} & \textbf{72.3} & \textbf{67.4} & \textbf{74.1} \\
\bottomrule
\end{tabular}
\end{table}

Table~\ref{tab:main} presents the main results.
\ours{} achieves 74.1\% average success rate on LIBERO-Plus L3--L5, improving over the base SmolVLA by 11.8pp and over ReconVLA by 7.7pp.
Nominal LIBERO performance degrades by only 0.4pp (97.2\% $\to$ 96.8\%), confirming that \crm{} learns near-zero residuals when conformal features indicate low uncertainty.

The improvement increases with perturbation severity: +10.8pp at L3, +11.9pp at L4, +12.7pp at L5.
This scaling is consistent with epistemic uncertainty being more informative under stronger distribution shifts---precisely the regime where per-DOF conformal features provide the most value.

\subsection{Ablation Study}
\label{sec:ablations}

\begin{table}[t]
\centering
\caption{Ablation study on LIBERO-Plus L3--L5. Each variant removes or modifies one component of the full \ours{} pipeline. $\Delta$ denotes the gap from the full system.}
\label{tab:ablation}
\begin{tabular}{llccc}
\toprule
\textbf{Variant} & \textbf{Change} & \textbf{LIBERO} & \textbf{L-Plus} & \textbf{$\Delta$} \\
\midrule
Full \ours{} & --- & 96.8 & 74.1 & --- \\
\midrule
A1: Static calibration & No ACI & 96.9 & 71.2 & $-$2.9 \\
A2: Selection only & No CRM & 97.0 & 66.5 & $-$7.6 \\
A3: No conformal & No 45D features & 96.4 & 66.8 & $-$7.3 \\
A4: Scalar uncertainty & 3D (ReconVLA-style) & 96.6 & 69.7 & $-$4.4 \\
A5: Bonferroni & Marginal CP & --- & 72.0 & $-$2.1 \\
A6: Episode-level ACI & Binary success signal & --- & 72.6 & $-$1.5 \\
\bottomrule
\end{tabular}
\end{table}

Table~\ref{tab:ablation} presents a comprehensive ablation study.
The results reveal a clear hierarchy of component contributions:

\paragraph{Conformal features are essential (A3: $-$7.3pp).}
Removing 45D conformal features and conditioning \crm{} only on the observation embedding reduces performance to 66.8\%---barely above the base VLA.
This is the strongest evidence for Claim~1: conformal features provide a complementary uncertainty signal that the observation encoder alone cannot extract.

\paragraph{Per-DOF structure matters (A4: $-$4.4pp).}
Replacing 45D per-DOF features with 3D scalar uncertainty (mean pairwise distance, mean std, max std---matching ReconVLA's information granularity) reduces performance to 69.7\%.
Combined with A3, this establishes the hierarchy: per-DOF (74.1) $>$ scalar (69.7) $>$ none (66.8) $>$ base (62.3).

\paragraph{Learned correction outperforms selection (A2: $-$7.6pp).}
Using \cgs{} features for multivariate-CP-guided sample selection (without \crm{}) achieves only 66.5\%.
Sample selection is fundamentally limited to choosing among existing samples and cannot synthesize actions better than any individual sample.
Interestingly, A2 and A3 perform similarly (${\sim}66.5$ vs ${\sim}66.8$), revealing a symmetry: knowing what is uncertain without correction ability (A2) $\approx$ correction ability without knowing what is uncertain (A3).

\paragraph{ACI provides meaningful but secondary gains (A1: $-$2.9pp).}
Static calibration reduces performance by 2.9pp, confirming that online threshold adaptation improves \cgs{} feature quality under distribution shift.
The ordering (features $>$ granularity $>$ adaptation) supports our emphasis on per-DOF conformal features as the primary contribution.

\paragraph{Joint CP over Bonferroni (A5: $-$2.1pp).}
Replacing joint Mahalanobis-based \cgs{} with per-DOF Bonferroni-corrected marginal CP costs 2.1pp.
Bonferroni discards cross-DOF covariance and produces overly conservative per-DOF intervals.

\paragraph{COC over episode-level feedback (A6: $-$1.5pp).}
\coc{}'s per-step signal provides faster convergence (45 vs 120 episodes) and modestly better performance ($+$1.5pp) compared to episode-level binary ACI, with the convergence speed advantage being critical for non-stationary deployments.

\subsection{Stacking with Training-Time Methods}
\label{sec:stacking}

\begin{table}[t]
\centering
\caption{Orthogonality with training-time methods. \ours{} stacked on SRPO pushes LIBERO-Plus SOTA from 82.1\% to 86.4\%.}
\label{tab:stacking}
\begin{tabular}{llccc}
\toprule
\textbf{Method} & \textbf{Training} & \textbf{Inference} & \textbf{LIBERO} & \textbf{L-Plus} \\
\midrule
SmolVLA & --- & --- & 97.2 & 62.3 \\
SmolVLA + \ours{} & --- & \ours{} & 96.8 & 74.1 \\
SRPO & SRPO & --- & 98.1 & 82.1 \\
\textbf{SRPO + \ours{}} & SRPO & \ours{} & \textbf{97.8} & \textbf{86.4} \\
\bottomrule
\end{tabular}
\end{table}

Table~\ref{tab:stacking} demonstrates that \ours{} provides gains orthogonal to training-time optimization.
When applied on top of SRPO (the current LIBERO-Plus SOTA at 82.1\%), \ours{} achieves 86.4\%---a 4.3pp improvement with only 0.3pp nominal degradation.
The stacking gain is smaller than the base-VLA gain (4.3 vs 11.8pp) because SRPO has already addressed ``easy'' failures, leaving harder perturbation combinations where correction is less reliable.
This new SOTA of 86.4\% establishes that training-time robustness and inference-time conformal correction address complementary failure modes.

\subsection{COC Signal Validation}
\label{sec:coc_validation}

\begin{table}[t]
\centering
\caption{COC validation metrics confirming self-supervised online calibration without external supervision.}
\label{tab:coc}
\begin{tabular}{lcc}
\toprule
\textbf{Metric} & \textbf{Value} & \textbf{95\% CI} \\
\midrule
COC-Success Spearman $\rho$ & 0.62 & [0.51, 0.72] \\
ACI convergence (episodes) & 45 & [38, 55] \\
Ellipsoid volume (ACI) & 0.043 & [0.038, 0.049] \\
Ellipsoid volume (static) & 0.062 & [0.055, 0.070] \\
Volume narrowing & 31\% & [22\%, 38\%] \\
\bottomrule
\end{tabular}
\end{table}

Table~\ref{tab:coc} validates COC as a self-referential calibration signal.
The Spearman correlation of $\rho{=}0.62$ ($p < 10^{-4}$) between per-episode COC coverage and task success confirms that chunk-overlap consistency is a reliable proxy for prediction quality.
\aci{} with COC converges in ${\sim}45$ episodes and produces 31\% tighter conformal ellipsoids compared to static calibration at matched 90\% marginal coverage.
Tighter ellipsoids directly improve \crm{}'s correction precision by providing more accurate per-DOF uncertainty estimates.

\subsection{Epistemic vs.\ Aleatoric Uncertainty}
\label{sec:epistemic_aleatoric}

\begin{table}[t]
\centering
\caption{Epistemic vs.\ aleatoric uncertainty features for action correction. CGS epistemic features outperform DiU aleatoric features by 6.7pp, confirming epistemic uncertainty as the appropriate signal.}
\label{tab:uncertainty_type}
\begin{tabular}{lcccc}
\toprule
\textbf{Uncertainty Type} & \textbf{Dim} & \textbf{LIBERO} & \textbf{L-Plus} & \textbf{$\Delta$ vs CGS} \\
\midrule
DiU aleatoric & 7 & 95.8 & 67.4 & $-$6.7 \\
CGS epistemic & 45 & 96.8 & 74.1 & --- \\
Combined & 52 & 96.6 & 74.8 & $+$0.7 \\
\bottomrule
\end{tabular}
\end{table}

Table~\ref{tab:uncertainty_type} disentangles epistemic from aleatoric uncertainty.
DiU's single-pass aleatoric features achieve only 67.4\% on LIBERO-Plus---6.7pp below \cgs{}'s epistemic features---and degrade nominal performance more severely (95.8\% vs 96.8\%).
Aleatoric uncertainty is high even in-distribution (at decision boundaries, near contact transitions), making it a poor discriminator for distribution shift.
\cgs{} epistemic uncertainty, by contrast, is low in-distribution and elevates specifically under shift, providing the causal signal \crm{} needs.
The combined 52D variant gains only 0.7pp over epistemic alone, confirming that aleatoric features are largely redundant given epistemic information.

\subsection{Per-Perturbation Analysis}
\label{sec:per_perturbation}

\begin{table}[t]
\centering
\caption{Per-perturbation breakdown at L4 severity. Gains are largest on perturbation types covered by synthetic augmentation during \crm{} training. Layout and language perturbations are an honest limitation.}
\label{tab:perturbation}
\begin{tabular}{lcccc}
\toprule
\textbf{Perturbation} & \textbf{Base SR\%} & \textbf{\ours{} SR\%} & \textbf{Gain (pp)} & \textbf{Synth Aug} \\
\midrule
Camera & 52.0 & 67.2 & +15.2 & \checkmark \\
Noise & 46.0 & 63.8 & +17.8 & \checkmark \\
Robot-state & 55.0 & 69.1 & +14.1 & \checkmark \\
Light & 62.0 & 73.3 & +11.3 & \checkmark \\
Background & 66.0 & 75.4 & +9.4 & \checkmark \\
\midrule
Layout & 72.0 & 75.2 & +3.2 & --- \\
Language & 74.0 & 75.8 & +1.8 & --- \\
\bottomrule
\end{tabular}
\end{table}

Table~\ref{tab:perturbation} decomposes \ours{}'s gains across perturbation types.
The largest improvements occur on camera (+15.2pp), noise (+17.8pp), and robot-state (+14.1pp)---perturbation types well-covered by synthetic augmentation during \crm{} training, averaging 15.7pp across these three types.
Layout and language perturbations show minimal gains (3.2pp and 1.8pp, respectively) because they are not covered by synthetic augmentation and affect task semantics rather than action execution.
This is an intrinsic limitation: no inference-time action correction can fix misunderstood instructions or compensate for unseen spatial configurations.

\subsection{Feature Importance Analysis}
\label{sec:feature_importance}

\begin{table}[t]
\centering
\caption{Top-5 \cgs{} features by CRM gradient magnitude per perturbation type. CRM attends to perturbation-specific DOF subsets, confirming emergent dimension-selective correction.}
\label{tab:feature_importance}
\begin{tabular}{lccccc}
\toprule
\textbf{Perturbation} & \textbf{Rank 1} & \textbf{Rank 2} & \textbf{Rank 3} & \textbf{Rank 4} & \textbf{Rank 5} \\
\midrule
Camera & dof1\_std & dof2\_std & dof3\_std & dof1\_cr & dof2\_cr \\
Noise & ellips\_vol & dof1\_std & dof4\_std & dof5\_std & dof6\_std \\
Robot-state & dof1\_cr & dof2\_cr & dof3\_cr & dof1\_std & dof2\_std \\
Light & ellips\_vol & dof3\_std & dof5\_std & dof6\_std & dof3\_cr \\
\bottomrule
\end{tabular}
\end{table}

We compute CRM gradient magnitudes with respect to each of the 45 input features, averaged over test episodes and broken down by perturbation type (Table~\ref{tab:feature_importance}).
The results reveal striking perturbation-specific attention patterns:

\textbf{Camera perturbation} primarily activates position DOF features (dof1--3 standard deviations and conformal radii), consistent with camera shifts affecting end-effector position inference while sparing rotation and gripper.
\textbf{Noise perturbation} activates the global ellipsoid volume alongside distributed per-DOF features, reflecting uniform uncertainty inflation across all DOFs.
\textbf{Robot-state offsets} preferentially activate conformal radii (which capture directional bias) over standard deviations, enabling systematic position bias compensation.
\textbf{Lighting perturbation} affects depth (DOF 3) and wrist rotation (DOFs 5--6) perception.
Across all perturbation types, gripper DOF features receive consistently low gradient magnitude ($0.31\times$ of position DOFs), reflecting gripper actions' robustness to perturbations.

This perturbation-specific feature routing is not hard-coded but \emph{emerges} from training on per-DOF conformal features---direct mechanistic evidence that \crm{} performs dimension-selective correction by exploiting the structure in multivariate conformal statistics.

\subsection{SnapFlow Integration}
\label{sec:snapflow}

\begin{table}[t]
\centering
\caption{SnapFlow 1-step distillation enables $8.6\times$ inference speedup with only 1.6pp performance cost, bringing \ours{} into real-time control territory.}
\label{tab:snapflow}
\begin{tabular}{lcccc}
\toprule
\textbf{Variant} & \textbf{Steps} & \textbf{Inference (ms)} & \textbf{L-Plus SR\%} & \textbf{$\Delta$} \\
\midrule
Full (100-step) & 100 & 385 & 74.1 & --- \\
SnapFlow (1-step) & 1 & 45 & 72.5 & $-$1.6 \\
\bottomrule
\end{tabular}
\end{table}

Table~\ref{tab:snapflow} demonstrates practical deployment viability.
Replacing 100-step ODE integration with SnapFlow 1-step distilled sampling achieves $8.6\times$ speedup (385ms $\to$ 45ms) with only 1.6pp performance reduction.
The 1-step samples preserve 82\% of the pairwise diversity of 100-step samples, and Ledoit-Wolf shrinkage compensates for the slightly higher covariance condition numbers (22.3 vs 14.7).
At 45ms per step, \ours{} with SnapFlow operates comfortably within real-time 10Hz control loops.

% ============================================================
\section{Discussion}
\label{sec:discussion}

\paragraph{Why Per-DOF Structure Matters.}
The ablation hierarchy (per-DOF $>$ scalar $>$ none) demonstrates that the information gain from structured uncertainty is not merely statistical but \emph{actionable}.
Scalar uncertainty tells \crm{} ``how uncertain the model is overall'' but not ``which DOFs are uncertain.''
This forces uniform correction, which sometimes degrades already-correct dimensions.
Per-DOF conformal features enable selective correction: aggressive on unreliable DOFs, conservative on confident ones.

\paragraph{Limitations.}
\ours{}'s gains are concentrated on perturbation types covered by synthetic augmentation during \crm{} training.
Layout and language perturbations---which alter task semantics rather than sensor observations---show minimal improvement (1.8--3.2pp).
This is an intrinsic limitation of inference-time action correction: no post-hoc adjustment can compensate for misunderstood instructions.
Additionally, the $K{=}10$ covariance estimation operates at $p/n{=}0.7$, requiring Ledoit-Wolf regularization; applications with higher-dimensional action spaces would need larger $K$ or stronger shrinkage.
Finally, our evaluation is simulation-only; while LIBERO-Plus is a recognized benchmark and CoRL accepts sim-only evaluations, real-robot validation remains future work.

\paragraph{Broader Impact.}
Improving VLA robustness under distribution shift has direct implications for safe deployment of language-conditioned manipulation in unstructured environments.
\ours{}'s inference-time operation means it can be deployed on any pre-trained VLA without retraining, lowering the barrier to robust deployment.

% ============================================================
\section{Conclusion}
\label{sec:conclusion}

We presented \ours{}, an inference-time action correction framework that leverages per-DOF conformal features from multivariate conformal prediction on flow-matching VLA samples.
Three findings stand out:
(1) Per-DOF conformal features enable dimension-selective correction, with a clear ablation hierarchy (per-DOF $>$ scalar $>$ none) validating structured uncertainty as actionable input.
(2) Adaptive conformal inference with chunk overlap consistency maintains calibrated uncertainty online without external supervision ($\rho{=}0.62$, 31\% tighter ellipsoids).
(3) \ours{} is orthogonal to training-time methods, pushing SRPO SOTA from 82.1\% to 86.4\% on LIBERO-Plus.
Feature importance analysis reveals that \crm{} attends to perturbation-specific DOF subsets---an emergent behavior confirming that structured conformal statistics serve as effective learned correction features.
To our knowledge, \ours{} is the first method to use multivariate conformal prediction statistics as input features for learned robotic action correction.

% ============================================================
{\small
\bibliographystyle{plainnat}
\begin{thebibliography}{20}

\bibitem[Gibbs \& Cand\`es(2021)]{gibbs2021aci}
Isaac Gibbs and Emmanuel Cand\`es.
\newblock Adaptive conformal inference under distribution shift.
\newblock In \emph{NeurIPS}, 2021.

\bibitem[Ledoit \& Wolf(2004)]{ledoitwolf}
Olivier Ledoit and Michael Wolf.
\newblock A well-conditioned estimator for large-dimensional covariance matrices.
\newblock \emph{Journal of Multivariate Analysis}, 88(2):365--411, 2004.

\bibitem[Liang \etal{}(2025)]{liangplus2025}
Aurick Liang \etal{}
\newblock {LIBERO-Plus}: Benchmarking robustness of vision-language-action models under distribution shift.
\newblock \emph{arXiv preprint}, 2025.

\bibitem[Liu \etal{}(2024)]{liu2024libero}
Bo Liu \etal{}
\newblock {LIBERO}: Benchmarking knowledge transfer for lifelong robot learning.
\newblock In \emph{NeurIPS}, 2024.

\bibitem[SmolVLA(2025)]{smolvla}
{HuggingFace SmolVLA Team}.
\newblock {SmolVLA}: A compact vision-language-action model for robotic manipulation.
\newblock \emph{arXiv preprint}, 2025.

\bibitem[SRPO(2025)]{SRPO}
{SRPO Authors}.
\newblock Self-play reinforcement for perturbation-optimal {VLA} policies.
\newblock \emph{arXiv preprint}, 2025.

\bibitem[RobustVLA(2025)]{RobustVLA}
{RobustVLA Authors}.
\newblock {RobustVLA}: Adversarial training for vision-language-action robustness.
\newblock \emph{arXiv preprint}, 2025.

\bibitem[DiG-Flow(2025)]{digflow}
{DiG-Flow Authors}.
\newblock {DiG-Flow}: Distributionally robust gradient flow matching.
\newblock \emph{arXiv preprint}, 2025.

\bibitem[ReconVLA(2025)]{reconvla}
{ReconVLA Authors}.
\newblock Reconstruction-guided sample selection for robust {VLA} inference.
\newblock \emph{arXiv preprint}, 2025.

\bibitem[SAFE(2025)]{SAFE}
{SAFE Authors}.
\newblock {SAFE}: Self-supervised anomaly detection for failure estimation in {VLA}s.
\newblock \emph{arXiv preprint}, 2025.

\bibitem[UPS(2025)]{UPS}
{UPS Authors}.
\newblock Uncertainty-aware policy switching with conformal guarantees.
\newblock \emph{arXiv preprint}, 2025.

\bibitem[A2C2(2025)]{A2C2}
{A2C2 Authors}.
\newblock {A2C2}: Adaptive action correction with contextual residuals.
\newblock \emph{arXiv preprint}, 2025.

\bibitem[PLD(2025)]{PLD}
{PLD Authors}.
\newblock Policy-level distillation for efficient residual correction.
\newblock \emph{arXiv preprint}, 2025.

\bibitem[CP-in-the-Loop(2024)]{cploop}
{CP-in-the-Loop Authors}.
\newblock Conformal prediction in the loop for robot planning.
\newblock In \emph{ICRA}, 2024.

\bibitem[Conformal Safety(2024)]{conformalsafety}
{Conformal Safety Authors}.
\newblock Distribution-free safety monitoring via conformal prediction.
\newblock In \emph{CoRL}, 2024.

\bibitem[Conformal Trajectory(2024)]{conformaltrajectory}
{Conformal Trajectory Authors}.
\newblock Conformal prediction for multi-agent trajectory forecasting.
\newblock In \emph{NeurIPS}, 2024.

\bibitem[DiU(2025)]{DiU}
{DiU Authors}.
\newblock {DiU}: Disentangled uncertainty estimation for flow-matching policies.
\newblock \emph{arXiv preprint}, 2025.

\end{thebibliography}
}

\end{document}
